\documentclass[11pt,a4paper]{article}

\usepackage[spanish,es-nodecimaldot]{babel}
\usepackage{csquotes}

%   \usepackage[general,draft,generous]{velvetblue}

\usepackage[general,final,compact,portada]{velvetblue}
\usepackage{amssymb}
\usepackage{enumitem}
\usepackage{setspace}

\VelvetPortadaLogo{logo.png}
\VelvetPortadaUniversidad{Universidad de las Américas Puebla (UDLAP)}
\VelvetPortadaLicenciatura{Licenciatura en Ciencia de Datos}
\VelvetPortadaActividad{Actividad 8: Procesamiento de Lenguaje Natural}
\VelvetPortadaAutores{%
    Andrea Hernandez Huerta\\
    ID: 175523\\[1em]
    Heriberto Espino Montelongo\\
    ID: 175199
}
\VelvetPortadaProfesor{Profesor: Dr. Victor Giovanni Morales Murillo}
\VelvetPortadaCurso{Curso: Programación Concurrente}
\VelvetPortadaFecha{Fecha de entrega: 25 de septiembre de 2026}

\setlist[enumerate]{
    itemsep=0pt,
    parsep=0pt,
    topsep=0pt,
    partopsep=0pt
}

\crefname{section}{sección}{secciones}
\Crefname{section}{Sección}{Secciones}
\crefname{subsection}{subsección}{subsecciones}
\Crefname{subsection}{Subsección}{Subsecciones}
\crefname{figure}{figura}{figuras}
\Crefname{figure}{Figura}{Figuras}
\crefname{table}{tabla}{tablas}
\Crefname{table}{Tabla}{Tablas}
\crefname{equation}{ecuación}{ecuaciones}
\Crefname{equation}{Ecuación}{Ecuaciones}
\crefname{enumi}{inciso}{incisos}
\Crefname{enumi}{Inciso}{Incisos}

\doublespacing

\addbibresource{references.bib}
\ExecuteBibliographyOptions{url=true}


\begin{document}

\vspace*{.4cm}
\noindent
Mucha información se encuentra expresada en
\textit{lenguaje natural}: publicaciones en redes sociales, mensajes, documentos,
reseñas o transcripciones. Estos textos se consideran
principalmente \textit{datos no estructurados}, por lo que no pueden ser utilizados
directamente por la mayoría de los algoritmos. Aquí aparece el
\textbf{Procesamiento de Lenguaje Natural (PLN)}, que busca desarrollar
métodos para representar, analizar y generar lenguaje humano.

El texto se divide en \textit{tokens}; posteriormente,
un \textit{embedding} asigna a cada token un vector para representar relaciones, permitiento usar medidas de similitud, como la similitud de coseno.
% \begin{equation}
%     \label{eq:similitud-coseno}
%     \operatorname{sim}(\mathbf{x},\mathbf{y})
%     =
%     \frac{\mathbf{x}^{\top}\mathbf{y}}
%     {\|\mathbf{x}\|\,\|\mathbf{y}\|},
% \end{equation}
A partir de estas representaciones pueden realizarse,
entre otras, las siguientes tareas:
\begin{enumerate}[label=\arabic*., leftmargin=*]
    \item \textbf{Clasificación y análisis de sentimientos:} determinar el tema,
    intención o polaridad de un texto.
    \item \textbf{Sistemas de recomendación:} utilizar contenido mediante
    similitudes entre vectores, filtrado colaborativo mediante matrices
    usuario--ítem y factorización matricial, o combinaciones híbridas.
    \item \textbf{Generación de lenguaje:} predecir sucesivamente nuevos tokens
    condicionados por el contexto anterior.
\end{enumerate}

Una conexión importante entre PLN y programación concurrente
aparece en los \textit{Transformers}. Antes de éstos, el procesamiento dependía
de \textit{RNNs}, que pueden presentar problemas de
\textit{vanishing} y \textit{exploding gradients}. Arquitecturas \textit{seq2seq}, mediante \textit{LSTMs}
y mecanismos anteriores de atención ayudaron a resolver parte de estas
limitaciones. \textcite{vaswani2017attention} propusieron el
Transformer, eliminando la recurrencia y permitiendo paralelismo en
distintas posiciones de una secuencia, calculando tres proyecciones
denominadas \textit{queries}, \textit{keys} y \textit{values}. Los Transformers actuales pueden
organizarse como modelos \textit{encoder--decoder}, como el trabajo original;
\textit{encoder-only}; o \textit{decoder-only}, familia a la que pertenece Chat 
\parencite{statquestNN,sandersonNN}.

Desde la \textbf{lingüística computacional}, el lenguaje también puede estudiarse
en distintos niveles:

\begin{enumerate}[label=\arabic*., leftmargin=*]
    \item \textbf{Fonética y fonología:} estudian los sonidos organzación dentro de una lengua, para reconocimiento de voz.
    \item \textbf{Morfología:} estudia la estructura interna de las palabras.
    Permite, por ejemplo, la \textit{lematización},
    $\textit{cantábamos}\rightarrow\textit{cantar}$, o el \textit{stemming},
    $\{\textit{trabajando},\textit{trabajadores}\}\rightarrow\textit{trabaj}$.
    \item \textbf{Sintaxis:} estudia cómo se relacionan las palabras dentro de una
    oración. ``\textit{Mary saw the man with a telescope}''.
    \item \textbf{Semántica y pragmática:} estudian el significado y su dependencia
    del contexto. ``\textit{Fui al Banco}''.
\end{enumerate}

La calidad de los datos es súper importante. Si las transcripciones o
etiquetas utilizadas durante el entrenamiento son incorrectas, el modelo
optimizará su función de pérdida respecto de información incorrecta.
Este problema se intensifica cuando existen pocos datos. 
% para las cuales existen pocos corpus digitalizados o datos anotados. 

% En México
% esto es particularmente relevante para lenguas indígenas como náhuatl, maya,
% zapoteco o totonaco. 
Por ello, la importancia del PLN en México es doble. Es el país con mayor hablantes que hablan español, la cuarta lengua más hablada; y al mismo tiempo, tiene una enorme  
diversidad lingüística, con lenguas indígenas como náhuatl, maya,
zapoteco o totonaco. Lo que  
plantea la necesidad de construir corpus,
anotaciones y herramientas computacionales para lenguas que actualmente cuentan
con muchos menos datos.

\printbibliography[title={Referencias}]

\end{document}
